% Generated draft: requires mathematical and source review.
\documentclass[11pt]{article}
\usepackage[margin=25mm]{geometry}
\usepackage{fontspec}
\setmainfont{Latin Modern Roman}
\usepackage{amsmath,amssymb,mathtools,unicode-math}
\setmathfont{Latin Modern Math}
\usepackage{enumitem,needspace,xurl,hyperref}
\setlength{\parindent}{0pt}
\setlength{\parskip}{5pt}
\setlength{\emergencystretch}{3em}
\newcommand{\N}{\mathbb{N}}
\newcommand{\Z}{\mathbb{Z}}
\newcommand{\Q}{\mathbb{Q}}
\newcommand{\R}{\mathbb{R}}
\newcommand{\C}{\mathbb{C}}
\newcommand{\F}{\mathbb{F}}
\newcommand{\A}{\mathbb{A}}
\newcommand{\PP}{\mathbb{P}}
\newcommand{\E}{\mathbb{E}}
\newcommand{\eps}{\varepsilon}
\newcommand{\dd}{\,\mathrm{d}}
\begin{document}
\section*{UnsolvedMath ID: 30000796}
\textbf{Source problem code:} catalogue entry
\subsection{Definitions and mathematical statement}
Let \(F\) be a field, let \(T\) be transcendental over \(F\), and let \(E\) be an elliptic curve over \(F(T)\). Write
\[
r_F(E)=\dim_{\Q}\bigl(E(F(T))\otimes_{\Z}\Q\bigr)
\]
for its rank, allowing this definition even when finite generation has not been established. In the characteristic-zero part of the question, the source restricts to ``nonconstant'' curves. A customary intended meaning is nonisotriviality, expressed by \(j(E)\notin F\); merely not having a model over \(F\) is a different condition, so the intended meaning must be checked.

The wording ``as \(E\) and/or \(F\) vary'' leaves two questions to distinguish. For a fixed characteristic-zero field \(F\), is there an integer \(B_F\geq0\) such that
\[
\forall E/F(T)\text{ in the intended nonconstant class},\qquad r_F(E)\leq B_F?
\]
Is there moreover one integer \(B\geq0\), independent of both \(F\) and \(E\), such that
\[
\forall F\ (\operatorname{char}F=0)\quad
\forall E/F(T)\text{ in that class},\qquad r_F(E)\leq B?
\]
The source singles out fields \(F\subseteq\C\), including the case \(F=\C\); this is a subclass of characteristic-zero fields, not a claim that every such field embeds in \(\C\). It also asks, more broadly, what can be said about ranks over \(F(T)\) for general fields \(F\), without giving a definite bound or specifying the intended scope in positive characteristic.

The raw statement calls \(E(F(T))\) finitely generated, whereas the cleaned statement omits that description. Finite generation and the treatment of constant or isotrivial curves and their constant-field contribution must be reviewed under the precise hypotheses; they are not imposed here without qualification.
\subsection{Short English statement}
For nonconstant elliptic curves over \(F(T)\) in characteristic zero, are the ranks uniformly bounded for fixed \(F\), or even by one bound as \(F\) also varies? The source also asks about general constant fields, without a precise universal claim.
\subsection{Sources}
\begin{itemize}
\item Local UnsolvedMath registry, permanent ID \(30000796\), source code catalogue entry, ``Uniform Rank Bounds for Elliptic Function Fields.'' No individual proposer is supplied.
\item \emph{Explicit Methods in Number Theory} (2008), Oberwolfach Reports, DOI 10.4171/owr/2007/34, pp. 1957--2034: \url{https://doi.org/10.4171/owr/2007/34}.
\item Supplied literature source, ``current elliptic-rank survey'': \url{https://www.math.uci.edu/~asilverb/bibliography/RanksCheatSheet.pdf}.
\item The literature triage checked 2026-08-21 labels characteristic-zero uniform rank boundedness unresolved, including C(T), and describes that case as already open in the report. This is a dated supplied assessment, not an independent literature verification.
\end{itemize}
\end{document}
